% ECONOMICS_PROBLEM: {"id": "E63-137", "title": "Joint monetary–fiscal policy under high debt and supply shocks", "jel_code": "E63", "source_id": "OP-0019"}
% !TeX program = xelatex
\documentclass[11pt,a4paper]{article}
\usepackage[margin=23mm]{geometry}
\usepackage{fontspec}
\setmainfont{Latin Modern Roman}
\usepackage{amsmath,amssymb,mathtools}
\usepackage{xcolor,enumitem,booktabs,longtable,array,xurl,hyperref}
\definecolor{muted}{HTML}{53636C}
\hypersetup{colorlinks=true,urlcolor=blue,linkcolor=blue}
\setlength{\parindent}{0pt}
\setlength{\parskip}{5pt}
\newcommand{\R}{\mathbb R}
\newcommand{\E}{\mathbb E}
\newcommand{\N}{\mathbb N}
\newcommand{\ind}{\mathbf 1}
\newcommand{\Var}{\operatorname{Var}}
\newcommand{\Cov}{\operatorname{Cov}}
\newcommand{\supp}{\operatorname{supp}}
\DeclareMathOperator*{\argmin}{arg\,min}
\DeclareMathOperator*{\argmax}{arg\,max}
\newcommand{\entrymeta}[3]{{\sffamily\small\color{muted}\textbf{\texttt{#1}}\quad|\quad #2\par #3\par}}
\newcommand{\Problemref}[1]{\texttt{#1}}
\newcommand{\JELref}[2]{\texttt{#2} (JEL \texttt{#1})}
\begin{document}
\section*{Joint monetary–fiscal policy under high debt and supply shocks}
\textbf{Problem ID:} \texttt{E63-137}\quad \textbf{JEL:} \texttt{E63}\par
% BEGIN ECONOMICS STATEMENT
\label{problem:OP-0019}
\label{atlas:prob:A9}

\noindent\textbf{Original atlas ID:} A9. \textbf{Formulation:} editorial specification of a research family.

\subsection*{Definitions and assumptions}

Let a specified stochastic general-equilibrium model have positive prices $P_t>0$, inflation $\pi_t=P_t/P_{t-1}-1$, activity gap $x_t$, real public debt $b_t=B_t/P_t$, and banking state $k_t$. A joint policy $p$ chooses nominal rate $i_t$, real purchases $G_t$, real net revenue $T_t$, and declared prudential instruments as functions of available histories. One-period debt issued at par obeys
\[
b_t=\frac{1+i_{t-1}}{1+\pi_t}b_{t-1}+G_t-T_t.
\]
Feasible policies satisfy $i_t\ge\underline i$, stated tax and spending limits, bank solvency/liquidity constraints, and the model's intertemporal government-budget and transversality conditions. Let $\mathcal E(p,\gamma)$ be equilibrium allocations under structural parameter $\gamma$; feasible policies must admit a nonempty such set for every admissible parameter. The welfare loss $L_\gamma(e)$ is specified from household welfare or an explicitly acknowledged quadratic approximation.

\subsection*{Formal research question}

For a declared nonempty identified parameter set $\Gamma_I$ and feasible policy class $\mathfrak P_F$, analyze
\[
\inf_{p\in\mathfrak P_F}\sup_{\gamma\in\Gamma_I}
 \sup_{e\in\mathcal E(p,\gamma)}L_\gamma(e),
\]
where the inner supremum makes equilibrium selection explicit. Characterize when commitments jointly stabilize inflation and activity while satisfying debt and banking constraints after a declared supply disturbance. Determine conditions for existence, equilibrium determinacy, and computational approximation. If monetary, fiscal, and prudential authorities have distinct objectives and choose strategically, replace the coordinated policy program with a specified policy game and its equilibrium concept. Compare that outcome with the coordinated optimum rather than treating the two as identical.

\subsection*{Interpretation and scope}

High debt alone does not define fiscal dominance. The response of fiscal backing to debt, the monetary reaction rule, and beliefs about future regimes determine the relevant equilibrium properties. A budget identity is necessary but insufficient for solvency; the intertemporal restriction and admissibility of price paths must also be stated. Lower-bound, banking, and debt limits can make local approximations misleading when large shocks cross constraints. A loss function weighting inflation, activity, and financial distress must be normatively justified and cannot be inferred from the words ``stabilizes the economy.'' Robust policy design also requires a credible ambiguity set, as unrestricted adverse models can make every policy look unacceptable.

\subsection*{Source audit and status}

Leeper (1991), Eggertsson--Woodford (2003), and Woodford (2003) are verified. They establish policy-regime determinacy and optimal monetary-policy results in specific models, with lower-bound analysis and welfare foundations. They do not jointly resolve all sovereign-debt, banking, supply-shock, and institutional constraints listed in the atlas. Conversely, the cited lower-bound optimal-policy benchmark is already solved and should not be relabeled open. The constrained minimax program above is editorial and may contain solved subcases. The substantive frontier is model-specific coordination, nonlinear constraints, and empirically supported policy comparisons, with an explicit distinction between common objectives and separate authorities.

\subsection*{References}

\begin{enumerate}

\item[S1] Eric M. Leeper (1991). \emph{Equilibria under `Active' and `Passive' Monetary and Fiscal Policies}. Journal of Monetary Economics 27(1), 129--147. \url{https://doi.org/10.1016/0304-3932(91)90007-B}. \textit{Locator:} article; policy-regime determinacy analysis. \textit{Establishes:} Dependence of equilibrium determination on monetary and fiscal responses. \textit{Does not establish:} a high-debt universal optimal-policy rule.

\item[S2] Gauti B. Eggertsson and Michael Woodford (2003). \emph{The Zero Bound on Interest Rates and Optimal Monetary Policy}. Brookings Papers on Economic Activity 2003(1), 139--235. \url{https://www.brookings.edu/articles/the-zero-bound-on-interest-rates-and-optimal-monetary-policy/}. \textit{Locator:} abstract; lower-bound policy analysis. \textit{Establishes:} Optimal history-dependent monetary policy in a specified lower-bound model. \textit{Does not establish:} optimal policy for all banking and sovereign-debt constraints.

\item[S3] Michael Woodford (2003). \emph{Interest and Prices: Foundations of a Theory of Monetary Policy}. Princeton University Press. \url{https://www.jstor.org/stable/j.ctv30pnvmf}. \textit{Locator:} publisher description; policy and welfare chapters. \textit{Establishes:} Foundations for interest-rate rules, nominal rigidities, and welfare-based policy. \textit{Does not establish:} a completed theory of every monetary--fiscal--banking interaction.

\end{enumerate}

\subsection*{Common conventions: Source mathematical conventions}

Unless an entry states otherwise, agent and firm sets are finite. State and action spaces are measurable, strategies and outcome maps are measurable, probability laws are specified, and expectations exist for the targets under discussion. Infinite-horizon discount factors lie in $(0,1)$ and discounted objectives are finite. Norms, tolerances and complexity input sizes must be specified for computational comparisons. Constants or primitives that appear locally are inputs, not empirical facts asserted by this catalogue.

\textbf{Identification.} A statistical model $m\in\mathcal M$ implies an observable law $P_m$ and target $\theta(m)$. At law $P$, its identified set is
\[
\Theta(P;\mathcal M)=\{\theta(m):m\in\mathcal M,\ P_m=P\}.
\]
Point identification holds when this set is a singleton, or equivalently when $P_{m_1}=P_{m_2}$ implies $\theta(m_1)=\theta(m_2)$ for all compatible models. Sharp bounds mean bounds attained, or approached when the boundary is not attained, by compatible models. Writing this set is a definition, not a solution: the substantive task is to establish informative restrictions and derive or estimate the set. An unrestricted-confounding model may give only trivial bounds.

\textbf{Inference.} Identification, estimability and valid uncertainty quantification are separate questions. Fix $\alpha\in(0,1)$. A confidence procedure $C_n$ over a declared law class $\mathcal P$ has asymptotic uniform coverage at level $1-\alpha$ when
\[
\liminf_{n\to\infty}\inf_{P\in\mathcal P}\Pr_P\{\theta(P)\in C_n\}\geq1-\alpha.
\]
For set-identified targets the coverage event must be defined separately, for example coverage of the full identified set. If an entry uses identification-indexed models instead of uniquely indexed laws, the same coverage convention is applied over those models. Pointwise limits do not establish uniform validity, and asymptotic validity does not guarantee finite-sample precision. Sample sizes, dependence structures and weak-identification sequences are part of each application.

\textbf{Counterfactuals.} $Y_i(a)$ is a potential outcome under a specified intervention, including its timing, implementation and versions. It is not $Y_i$ conditional on voluntarily choosing $a$. A full intervention vector permits interference; shorthand scalar treatment notation does not silently impose no interference. Assignment, selection, attrition, anticipation and measurement must be specified. A claim about an instrument's local effect is not automatically a population policy effect.

\textbf{Economic equilibrium.} A counterfactual model includes best responses, constraints, feasibility, market clearing, state transitions and expectations consistent with the stated information. When several equilibria exist, either a declared selector is an input or the target is set-valued. Comparisons across policy regimes must use the stated selection convention. Reduced-form relationships are not presumed invariant after an equilibrium-changing intervention.

\textbf{Optimization and welfare.} Optimization is over the stated feasible set. If attainment is not established, $\sup$ or $\inf$ and $\varepsilon$-optimal choices replace an asserted maximizer or minimizer. For example, with value $V=\sup_{a\in\mathcal A}W(a)<\infty$, an $\varepsilon$-optimal set is $\{a\in\mathcal A:W(a)\geq V-\varepsilon\}$ for $\varepsilon>0$. Benefits, costs, risk and health or environmental outcomes can be aggregated only using declared units and conversion weights. Interpersonal utility weights, the treatment of fiscal transfers and foreign profits, discounting, and the behavioral welfare criterion are explicit inputs. A comparative-static derivative additionally requires existence and appropriate differentiability of the selected counterfactual.

\textbf{Sources.} Local labels [S1], [S2], and so forth restart in every question. Locators identify the relevant abstract, model, section or result at the level actually checked; a bibliographic or abstract-level verification is not represented as inspection of an entire proof. Working papers are distinguished from later publications and changing versions. Source summaries are paraphrased. All URL checks and editorial formulations refer to the audit date stated above.
% END ECONOMICS STATEMENT
\end{document}
